% =============================================================================
% 4.1 与 4.2 节图表索引（已分发，本文件仅留作清单参考）
%
% 版本：v0.2 (2026-05-05)
%   v0.1：作为图表汇总文件，集中存放 6 张表的 LaTeX 源码。
%   v0.2：6 张表已分发到 4_1_语义桥接.tex 与 4_2_代码生成.tex 中对应的
%         \reftab 引用段附近，便于浮动机制就近放置；2 张图同步嵌入。
%         本文件不再 \input 到主稿，仅作为图表清单查阅。
%
% 注意：请勿将本文件 \input 到论文主稿，否则 \label 会发生重复定义。
% =============================================================================


% --------------------------------------------------------
% 4.1 节图表清单（已嵌入 4_1_语义桥接.tex）
% --------------------------------------------------------
% 图 4-1  fig:dual-path-rag       4.1.1 节末段，图片 figures/dual_path_rag
% 表 4-1  tab:rag-weight-sweep    4.1.2 节，60 条 × 8 个权重点
% 表 4-2  tab:semantic-mapping    4.1.3 节，5 类要素分级阈值映射节选 19 条
% 表 4-3  tab:bridge-ablation     4.1.3 节，通路 B 4 档消融总表（71 条）

% --------------------------------------------------------
% 4.2 节图表清单（已嵌入 4_2_代码生成.tex）
% --------------------------------------------------------
% 图 4-2  fig:code-gen-flow       4.2.1 节首段，图片 figures/code-gen-flow
% 表 4-4  tab:code-bench-design   4.2.2 节，18 条评测集任务分布
% 表 4-5  tab:code-eval-overall   4.2.2 节，3 档总表
% 表 4-6  tab:code-eval-category  4.2.2 节，按类别拆分


% --------------------------------------------------------
% 图片接入说明（详见 figures/README.md）
% --------------------------------------------------------
% xelatex 不直接支持 .svg 文件。两种推荐做法：
%   方案 1（推荐）：在 drawio 中文件 → 导出 → PDF，存为同名 .pdf 与 .svg
%                  并列。\includegraphics 不写后缀，xelatex 自动选 .pdf。
%   方案 2（备选）：在导言区追加 \usepackage{svg}，编译时启用
%                  -shell-escape 并安装 inkscape，使用 \includesvg{...}。
% 当前 4.1 与 4.2 草稿用的是方案 1，引用 figures/dual_path_rag 与
% figures/code-gen-flow，不带后缀。
